%%=============================================================================
%% Voorwoord
%%=============================================================================

\chapter*{\IfLanguageName{dutch}{Woord vooraf}{Preface}}%
\label{ch:voorwoord}

Ik speel al meer dan twaalf jaar Kerbal Space Program. Dat klinkt misschien overdreven, maar wie het spel kent, begrijpt het meteen: KSP is geen gewoon spel. Het is een ruimtevaartsimulator waar je zelf raketten bouwt, missies plant en (als alles goed gaat) succesvol landt op andere planeten. Voor iemand die geïnteresseerd is in raketten en ruimtevaart, is het eigenlijk het perfecte spel.

Mijn primaire KSP installatie is met het echte zonnestelsel nagebootst in het spel, wat een modpack van 130 mods vereist. En wie mods installeert, weet ook wat er vroeg of laat fout gaat: mysterieuze crashes, mods die elkaar tegenwerken, en uren spenderen aan debuggen om uiteindelijk toch opnieuw te beginnen. Dat heb ik in die twaalf jaar meer keren meegemaakt dan ik wil toegeven. Dit probleem is dan ook precies wat ik in deze bachelorproef probeer aan te pakken. Het voelde meteen als een onderwerp waar ik écht iets mee kon doen, niet alleen voor mezelf, maar ook voor de bredere KSP-community.

Graag bedank ik een aantal mensen die dit mogelijk hebben gemaakt. Harvester en Squad verdienen een grote dank voor het creëren van dit unieke spel. De modding-community, en in het bijzonder BallisticFox voor al het werk in de nieuwe versie van de Real Solar System mod, houdt het spel levend lang nadat de officiële ontwikkeling stopte in 2023. Verder hebben YouTubers zoals Scott Manley, Matt Lowne, Everyday Astronaut, CSI Starbase en het NASASpaceflight-team mijn fascinatie voor de ruimte jaar na jaar verder aangewakkerd.

Een bijzonder grote dank aan mijn promotor, dhr.\ Andreas De Witte, en co-promotor, dhr.\ Sion Verschraege. Hun oprechte interesse in het onderwerp en hun waardevolle feedback hebben deze bachelorproef naar een hoger niveau gebracht.

Tot slot wil ik mijn familie en vrienden van harte bedanken. Thuis kon ik altijd rekenen op een luisterend oor als ik vast zat, en me er altijd bovenop hielpen als de stress te groot werd. Ik kon ook altijd rekenen op mijn vrienden voor de leuke game-sessies of gewoon altijd beschikbaar waren om me even af te leiden, jullie weten wie jullie zijn. Zonder die steun was dit een stuk zwaarder geweest.

\bigskip

\noindent Jules Van Nevel, 26 mei 2026